\cotefile{9}
\scanpage{16}
$A_0$ VA sur corps $k$ de car. $>0$.

Question : quand est-il possible, sur tte déf. inf. $A/\Lambda$ de $A_0$
(ou $\Lambda$ $k$-algèbre artinienne locale) d'avoir une
$k$-stratification sur $\mathrm{DR}^1(A/\Lambda)$, compatible avec
\unsure{changmt|changt} de base et fonctoriel en $A$ ?

Si $M$ est la \unsure{var.} des modules formels \deleted{\unsure{des|de}} $A_0$ sur $k$,
\deleted{il faut} et $A_M$ la déformation modulaire, il
faut que $\mathscr H=\mathrm{DR}^1(A_M/M)$ ait une \deleted{\unsure{connexion}}\inserted{stratification} formelle
\uline{compatible avec les opérations} $F,V:\mathscr H^{(p)}\mathrel{\substack{\xrightarrow{F}\\[-.5ex]\xleftarrow{V}}}\mathscr H$. Cela
implique que pour tt $x\in M$, le type (géométrique) de $\mathscr H_x$ muni de $F_x,V_x$ est le $=$ que
celui de $H_0=\mathrm{DR}^1(A_0)$, muni de $F_0,V_0$. Cela
indique sans doute que $A_0$ est ``ordinaire'' i.e.

\diagram{\begin{tikzcd}[ampersand replacement=\&,column sep=small,row sep=small]
{}\arrow[r,dash] \& H^{10}\arrow[r,dash]\arrow[d,phantom,"\deleted{\unsure{F}}"'] \& H\arrow[r,dash]\arrow[d,shift right,"V"'] \& H^{\unsure{01|0*}}\arrow[r]\arrow[d,phantom,"\deleted{\unsure{F}}"] \& {}\\
{}\arrow[r,dash] \& H^{10(p)}\arrow[r,dash] \& H^{(p)}\arrow[u,shift right,"F"']\arrow[r,dash] \& H^{01(p)} \& {}
\end{tikzcd}}
\[
\operatorname{Im}V=\operatorname{Ker}F=H^{10(p)},\quad
\operatorname{Ker}V=\operatorname{Im}F\unsure{'\,|,}\ ,
\]
\unsure{a|c'est} le nb max. de pts d'ordre $p$, ou ce qui revient
au $=$, que dans $t_{A_0^*}=H^1(A_0,\underline O_{A_0})
=H^{01}(A_0)$, la puissance p.ième
soit non dégénérée [i.e. induise un isom de
$H^{01(p)}$ sur $H^{\unsure{01|0*}}$]. // Mais alors \deleted{il y} pour tt $A$ sur $\Lambda$,
il en est de $=$ dans $\mathrm{DR}^1(A)$. Cela implique
que l'ext. $H=\mathrm{DR}^1(A)$ de $H^{01}$ par $H^{10}$ \unsure{splitte|splite|splitté|splittée} can.,
par $\operatorname{Im} F=F(H^{(p)})$, et \deleted{\unsure{le}} $\operatorname{Im} F\inserted{=H'}$ \unsure{est|$\longrightarrow$} évidemment
% archive sheet number: 15
\scanpage{17}
stable pour la stratification envisagée. De $=$, comme $VF=0$,
$V$ est nul sur $H'$, et \unsure{son} image étant contenue
dans $H^{10\,(p)}\inserted{=\operatorname{Ker}F}$\inserted{ (\unsure{donc $=$|étant $=$|équiv. à} $\operatorname{Ker}V=\operatorname{Im}F=H'$)}, \deleted{\unsure{$\operatorname{Ker}F(H^{(p)})$|car $FV=0$}},
et $=$ \unsure{égale|icelle} ; on trouve
que \uline{$H^{10\,(p)}$ est stable par \deleted{\unsure{connexion}}\inserted{la stratification}}. \deleted{donc}

\uline{Attention}, cela \uline{n'implique pas à priori que $H^{10}$ le
soit également}. Tout ce qu'on peut dire (valable
aussi pour $\Lambda$ une algèbre locale noeth. complète,
corps résiduel $k$ \ldots) c'est que si $S=\operatorname{Specf}(\Lambda)$,
et si on considère les deux images inverses
$E_1,E_2$ de $H$ sur \deleted{\unsure{$S$}} $T=\operatorname{Specf}(\Lambda\widehat\otimes\Lambda)=S\times S$,
\deleted{\unsure{alors}} et les deux images inverses $F_1,F_2$ de $H^{10}$,
et considérant l'isom de stratification
$u:E_1\xrightarrow{\sim}E_2$, d'où $u^{(p)}:E_1^{(p)}\mathbin{\unsure{-|\to}}E_2^{(p)}$,
alors $u^{(p)}(F_1^{(p)})=\deleted{\unsure{u^{(p)}}}\unsure{(|{}}F_2^{(p)}$, i.e.
$u(F_1)^{(p)}=F_2^{(p)}$. Si donc $u(F_1)=G$, $F_2=G'$, $E=E_2$
on aura $G^{(p)}=G'^{(p)}$ comme \deleted{\unsure{élémts}}\inserted{sections} de
de $\operatorname{Grass}(E)$ sur $T$, ou ce qui revient au
$=$, ces sections deviennent égales \deleted{par} \unsure{changmt|changt} de base
$\underset{\substack{\shortparallel\\T\ \text{puissance}\\\text{p.ième}}}{T'}\xrightarrow{\pi}T$. On ne peut en général en
conclure $G=G'$, car en général $\pi$ n'est pas un
épim. Mais si \inserted{$S$ donc} $T$ est formellt lisse sur $k$,
% archive sheet number: 16
\scanpage{18}
cela marchera ($\pi$ est alors \unsure{fidlt|fid.} plat). Or cela suffit ici,
en se ramenant au cas $S=M$ (variété des modules
formels). Donc $H^{10}$ \uline{est stable pour la stratification}.

Donc $H$ admet une décomposition fonctorielle
et stable par ext. de la base
\[
H\simeq H^{10}+H^{01}
\]
et cette décomposition est \uline{stable par stratification}.

Donc, le pb est maintenant de trouver une
stratification \deleted{\unsure{can.}} (fonctorielle ds les deux sens)
sur les deux morceaux $H^{10}=t_A^\vee$ et $H^{01}=t_{A^*}$,
Le deuxième pb, \inserted{pour $A_0$,} se ramène au premier pour $A_0^*$.

Donc il faut trouver une stratification fonctorielle sur
$t_A$. Or il y en a déjà \inserted{une}, grâce à la structure
de puissances p.ièmes non dégénérée, qui
\margin{à préciser \ldots}
% slanted
permet \unsure{de introduire|d'introduire|d'intr. d'emblée} le \unsure{grpe|groupe} structural
discret $\mathrm{Gl}(g)_k$. Et c'est la seule, \deleted{\unsure{(pour $\Lambda,A$ donnés|(pour $A,A$ donnés}} comme on voit
aisément en utilisant la compatibilité avec \unsure{la|les}
puissances p.ièmes (pour $\Lambda$, $A$ sur $\unsure{\lambda|\Lambda|X}$, donnés).

\uline{Remarque}\quad On voit donc que la position de
la filtration p.r. à la stratification ne peut aucunement
donner un principe de classification \inserted{des déf.} infinitésimale en car. $p>0$ !
% archive sheet number: 17
\scanpage{19}
La décomposition envisagée (can, fonct.)
\[
\mathrm{DR}^1(A/S)\simeq\underline{t}_A^\vee+\underline{t}_{A^*}
\]
vaut pour tt préschéma abélien \uline{ordinaire} sur
un $S$ de car. $p>0$, en $=$ temps que les expressions
\[
\underline{t}_A\simeq\unsure{\Phi|\mathscr L|\mathcal L}_A\otimes_{\mathbb F_p}\underline O_S,\qquad
\underline{t}_{A^*}\simeq\Phi_{A^*}\otimes_{\mathbb F_p}\underline O_S
\]
où $\Phi_A,\Phi_{A^*}$ sont des faisceaux \inserted{étales} loc. \deleted{\unsure{constants|const.}}
libres de rang $g$ ($=$ dim rel de $A/S$) d'esp vectoriels
sur $\mathbb F_p$. Si $S$ est le spectre d'un corps parfait,
on a
\[
\begin{aligned}
\Phi_{A^*}=(\underline t_{A^*})^F&=H^1(A,\underline O_A)^F=H^1(A,\mathbb Z/p\mathbb Z)\\
&=\operatorname{Hom}({}_pA\unsure{'|\cdot|{}},\mathbb Z/p\mathbb Z)=\deleted{\unsure{\Phi_A}}{}_pA(k)^\vee
\end{aligned}
\]
\[
\left\{\begin{aligned}
\Phi_{A^*}\deleted{\unsure{k}}(k)&={}_pA(k)^\vee\\
\Phi_A(k)&={}_pA^*(k)^\vee
\end{aligned}\right.
\]
Donc
\[
\left\{\begin{aligned}
\underline t_A^\vee&=\Phi_A\deleted{\unsure{(k)}}^\vee\otimes_{\mathbb F_p}k={}_p A^*(k)\otimes_{\mathbb F_p}k\\
t_{A^*}&={}_pA(k)^\vee\otimes_{\mathbb F_p}k .
\end{aligned}\right.
\]

On a des relations analogues pour $S$ (de car $p>0$)
quelconque : l'hyp. sur $A$ implique que
$\,{}_pA\mathbin{\to}S$ s'écrit de façon unique comme
composé $\,{}_pA\longrightarrow({}_pA)_{\mathrm{\acute et}}\mathbin{-}S$ de deux morphismes
finis loc libres, le premier radiciel surj, le $2^{\mathrm{ième}}$ \uline{étale}.
% archive sheet number: 18
\scanpage{20}
Ceci dit, on a des isom. can.
\[
\left\{\begin{aligned}
\unsure{\Phi|\mathscr L|\mathcal L}_A&=({}_pA^*)_{\mathrm{\acute et}}^\vee\\
\Phi_{A^*}&=({}_pA)_{\mathrm{\acute et}}^\vee
\end{aligned}\right.
\]

\uline{NB.} Soit $\varphi$ une polarisation de $A$ (plus
gén. une \deleted{\unsure{diviseur}} section de $\underline{\mathrm{NS}}_{A/S}$). Alors
$\varphi$ définit une \deleted{forme alternée sur}\inserted{section \deleted{\unsure{dite de}} $\varphi'$ de}
$\bigwedge^2\mathrm{DR}^1(A/S)$, \deleted{\unsure{à} valeurs en $\underline O_S$. Cette forme admet
les composantes $\underline H^{10}=\underline t_A^\vee$ et $H^{01}=\underline t_{A^*}$ comme
sous-espaces isotropes, et définit donc un
accouplement $\underline t_A^\vee\otimes\underline t_{A^*}\mathbin{\to}\underline O_S$ i.e.
un hom $\underline t_{A^*}\longrightarrow\underline t_A$. Cet hom n'est
autre que l'hom tangent à $\widetilde\varphi:A\mathbin{-}A^*$}
\[
\bigwedge^{2}\mathrm{DR}^1(A/S)\simeq\bigwedge^2\underline t_A^\vee+
\overbrace{\underline t_A^\vee\otimes\underline t_{A^*}}^{\operatorname{Hom}(\underline t_A,\underline t_{A^*})}
+\bigwedge^{\unsure{2}}\underline t_{A^*}
\]
Ceci dit, $\varphi'$ est dans le terme médian, et
comme tel n'est autre que l'hom
tangent $\widetilde\varphi':\underline t_A\mathbin{-}\underline t_{A^*}$ de $\widetilde\varphi:A\to A^*$ défini
par $\varphi$, qui est déterminé \unsure{également|ipso facto|formellement}, par tensorisation
\deleted{par} $\otimes_{\mathbb F_p}\underline O_S$, par
\diagram{\begin{tikzcd}[ampersand replacement=\&]
\underline\Phi_A \arrow[r,dash] \arrow[d,phantom,"\shortparallel" description] \& \underline\Phi_{A^*} \arrow[d,phantom,"\shortparallel" description]\\
({}_pA^*)_{\mathrm{\acute et}}^\vee \arrow[r,dash] \& ({}_pA)_{\mathrm{\acute et}}^\vee
\end{tikzcd}} transposé de
$({}_pA)_{\mathrm{\acute et}}\mathbin{-}({}_pA^*)_{\mathrm{\acute et}}$ induit par $\,{}_p\widetilde\varphi:{}_pA\mathbin{-}{}_pA^*$ \ldots
% archive sheet number: 19
\scanpage{21}
\uline{Re N.B.} Soit tjrs $g=$ dim rel. de $A/S$, et considérons
$\omega_A=\deleted{A}\operatorname{d\acute et}(\underline t_A)^{-1}
=\bigwedge^g\underline t_A\deleted{\unsure{{}^\vee}}$
\inserted{(formes diff de degré max le long de la section unité)}.
C'est un faisceau inv.
à p-structure non dégénérée, i.e. déduit par
ext. des scalaires d'un \deleted{\unsure{torseur}}\deleted{\unsure{sur}}\inserted{\deleted{\unsure{fais|faisc.}} faisceau loc. ct}
\deleted{\unsure{$\mathbb F_p^*$|$({}_pA^*)_{\mathrm{\acute et}}$}}\unsure{de} \unsure{inv|dim} sur $\mathbb F_p$, savoir
$\deleted{\unsure{(|\{}}\operatorname{d\acute et}\unsure{{}_g|{}_p|{}}({}_pA^*)_{\mathrm{\acute et}}\deleted{\unsure{)|\}}}$.
À ce titre, $\omega_A^{\otimes(p-1)}$ admet une section canonique.

Si on ne suppose plus \unsure{néc.} $A/S$ ordinaire, et si $U$ est l'ouvert des $s\in S$
tels que $A_s$ soit ordinaire, tt ce qui
précède s'applique à $A|U$, en particulier
on trouve une section canonique \inserted{$\sigma$} de
$L|U$, partout $\ne0$. \deleted{Si $U$ est sch. dense dans $S$,
on peut considérer $\sigma$ comme une \deleted{\unsure{fon}}\inserted{pseudo section}
\unsure{de $L$} sur $S$ (ou fonction méromorphe sur $S$), \deleted{\unsure{qui se}}\inserted{mais en}}

Mais en fait, \unsure{elle|on} provient d'une section
canonique de $L^{\otimes(p-1)}$ sur $S$ tt entier, défini
à l'aide de
\[
t_A^{(p)}\xrightarrow{\pi}t_A
\]
en passant au dét, d'où
\[
(\operatorname{d\acute et}\underline t_A)^{\otimes p}\longrightarrow\operatorname{d\acute et}\underline t_A
\]
d'où
\[
\sigma_A\in\Gamma(\operatorname{det}\underline t_A)^{1-p}=\Gamma\omega_A^{(p-1)}
\]
% archive sheet number: 20
\scanpage{22}
Cette section mérite le nom de \uline{invariant de
Hasse} de $A$. L'ens des zéros de cette
section est \unsure{exactmt|exactement} l'ens des $s$ tels que
$A_s$ soit \uline{non} \uline{ordinaire}.

Quelles sont les relations entre $\omega_A$ et $\omega_{A^*}$ ? On
a, grâce à la suite exacte
\[
\unsure{0|\cdot}\mathbin{\unsure{-|\to}}\underline t_A^\vee\to\mathrm{DR}^1(A/S)\mathbin{-}\underline t_{A^*}\mathbin{-}0
\]
la relation
\diagram{\begin{tikzcd}[ampersand replacement=\&,column sep=small,row sep=tiny]
\det t_A^\vee\otimes\det\underline t_{A^*}\arrow[r,phantom,"\simeq"] \& \det\mathrm{DR}^1(A/S)\arrow[d,phantom,"\unsure{\simeq|=|\sim}"{description,rotate=90}]\\
\& \mathrm{DR}^{2g}(A/S)\simeq\underline O_S
\end{tikzcd}}
donc on aura
\[
\omega_A\cdot\omega_{A^*}^\vee\simeq\underline O_S
\]
i.e. on a un isom can.
\[
\omega_A\simeq\omega_{A^*}
\]
Je dis que \uline{cet isom est compatible avec les
p-structures}, \deleted{\unsure{et|---}} \uline{donc transforme} $\sigma_A$ \uline{en} $\sigma_{A^*}$, \uline{et
dans le cas} $A$ \uline{ordinaire}, induit un isom
\[
\operatorname{d\acute et}({}_pA)_{\mathrm{\acute et}}\simeq\operatorname{d\acute et}({}_p A^*)_{\mathrm{\acute et}} .
\]
% archive sheet number: 21
\scanpage{23}
Il s'impose de regarder \unsure{égalmt|égalt}, si $A$ est ordinaire
\[
T_p(A)_{\mathrm{\acute et}}=\varprojlim\deleted{\unsure{T_p}}({}_{p^{\unsure{\nu|n}}}A)_{\mathrm{\acute et}}
\]
(N.B. $T_p(A)$ est une extension
\[
0\to T_p(A)_{\mathrm{rad}}\to T_p(A)\mathbin{\to}T_p(A)_{\mathrm{\acute et}}\to0
\]
On sait, par dualité de Cartier, que dans la
suite exacte correspondante
\[
0\mathbin{-}T_p(A^*)_{\mathrm{rad}}\mathbin{-}T_p(A^*)\mathbin{-}T_p(A^*)_{\mathrm{\acute et}}\mathbin{-}0
\]
les termes médians sont duals l'un de
l'autre, et \deleted{plus précisément} cette dualité est
compatible avec \unsure{filtration|filtrations}, donc
\[
\begin{aligned}
T_p(A)_{\mathrm{rad}}&\simeq D\bigl(T_p(A^*)_{\mathrm{\acute et}}\bigr)\\
T_p(A^*)_{\mathrm{rad}}&\simeq D\bigl(T_p(A)_{\mathrm{\acute et}}\bigr)
\end{aligned}
\]
donc la connaissance de $T_p(A)_{\mathrm{\acute et}}$ et $T_p(A^*)_{\mathrm{\acute et}}$
permet de reconstituer les \uline{deux} facteurs
de $T_p(A)$, donc $T_p(A)$ si la base est
parfaite [en général, il y aurait lieu d'
expliciter la structure d'extension \ldots].

On aurait envie d'établir un isom can.
\deleted{fonctoriel}
% archive sheet number: 22
\scanpage{24}
\[
\operatorname{d\acute et}T_p(A)_{\mathrm{\acute et}}\xrightarrow{\sim}\operatorname{d\acute et}T_p(A^*)
\]
[\deleted{Soit} Si on a $l\ne p$ au lieu de $p$, on a bien
\[
\det T_{\ell}(A)\overset{\alpha_\ell}{\simeq}\det T_\ell(A^*)\simeq\mathbb Z_\ell(g)\quad]
\]
\deleted{car $T_\ell(A^*)\simeq\deleted{\det}T_\ell(A^{*})^\vee\otimes\mathbb Z_\ell(1)$
donc $\det T_\ell(A^*)\simeq\det T_\ell(A)$
(Choisissons une polarisation)} Si $\varphi$ est une
polarisation de $A$,
définissant $\widetilde\varphi:A\mathbin{\unsure{-|\to}}A^*$ d'où $T_\ell(\widetilde\varphi):T_\ell(A)\mathbin{-}T_\ell(A^*)$,
$\operatorname{det}T_\ell(\widetilde\varphi):\operatorname{d\acute et}T_\ell(A)\xrightarrow{\beta_\ell(\varphi)}\operatorname{d\acute et}T_\ell(A^*)$,
alors l'isom $\alpha_\ell$ n'est autre que $\beta_\ell(\varphi)/\deg\varphi$. En s'inspirant de
cette formule, on définira,
\[
\beta_p(\varphi)=\operatorname{d\acute et}\bigl(T_p(\widetilde\varphi)_{\mathrm{\acute et}}\bigr):
\operatorname{d\acute et}T_p(A)_{\mathrm{\acute et}}\mathbin{-}\operatorname{d\acute et}T_p(A^*)_{\mathrm{\acute et}}
\]
et on considère
\[
\alpha_p(\varphi)=\beta_p(\varphi)/\sqrt{\deg\widetilde\varphi}:
\det T_p(A)_{\mathrm{\acute et}}\inserted{\otimes\unsure{\mathbb Q_p|\mathbb Q_\ell}}\mathbin{-}
\operatorname{d\acute et}T_p(A^*)_{\mathrm{\acute et}}\inserted{\otimes\mathbb Q_\ell}
\]

À prouver : a) $\alpha_p(\varphi)$ ne dépend pas de $\varphi$.

\deleted{\unsure{Cela|Ceci} \unsure{proviendrait|trivialement} formellt du \unsure{cas}}

\deleted{\uline{Lemme}\quad Soit $u$ un \deleted{\unsure{auto}}\deleted{\unsure{endom}} endom de $A$, d'où
\[
[\operatorname{d\acute et}T_p(u)_{\mathrm{\acute et}}:\operatorname{d\acute et}T_p(A)_{\mathrm{\acute et}}
\deleted{\unsure{\mathbb Z_p}}\to\operatorname{d\acute et}T_p(A)_{\mathrm{\acute et}}]\in\mathbb Z_p .
\]
On a}
\margin{\uline{faux}}
% slanted
\deleted{\[
\left\{\begin{aligned}
&\deg u\text{ est un carré parfait}\\
&\operatorname{d\acute et}\bigl(T_p(u)_{\mathrm{\acute et}}\bigr)=\sqrt{\deg u}\qquad\text{i.e.}\\
&\deg u=\bigl(\operatorname{d\acute et}(T_p(u)_{\mathrm{\acute et}})\bigr)^{2}\quad\text{\uline{et}, }\det T_p(u)_{\mathrm{\acute et}}\text{ est un entier }\unsure{>0|\geq0}
\end{aligned}\right.
\]}
% archive sheet number: 23
\scanpage{25}
En effet, soit $\psi$ une autre polarisation, donc
$\widetilde\psi=\widetilde\varphi\circ u$, où $u\in\operatorname{End}(A)\otimes_{\mathbb Z}\mathbb Q$, \deleted{\unsure{la}} et la condition
$\widetilde\psi'=-\widetilde\psi$ (\unsure{jointe à|joint à|pour :} $\widetilde\varphi'=-\widetilde\varphi$) donne
$\underset{\substack{\shortparallel\\-u'\widetilde\varphi}}{u'\widetilde\varphi'}=-\widetilde\varphi u$
i.e. $u'=\widetilde\varphi u\widetilde\varphi^{-1}$. On veut
prouver [posant $\delta(\widetilde\varphi)=\sqrt{\deg\widetilde\varphi}$, $\alpha_p(\varphi)=\Delta(\widetilde\varphi)$]
\[
\begin{aligned}
\Delta(\unsure{\widetilde\psi|\widetilde\varphi})/\delta(\unsure{\widetilde\varphi|\widetilde\psi})&=\Delta(\widetilde\varphi)/\delta(\widetilde\varphi)\\
\shortparallel\qquad&\\[-1ex]
\Delta(\widetilde\varphi)\circ\Delta(u)/\delta(\widetilde\varphi)\,\delta(u)&
\end{aligned}
\]
qui équivaut à : $\Delta(u)/\delta(u)=1\unsure{{}_{\ell}|{}}$. Or par
\unsure{Serre|dual.|lemme} on a $\deleted{\Delta}\Delta(u)\Delta(u^*)=\underset{\substack{\shortparallel\\\deg u}}{\delta(u)^2}$

d'autre part, comme $u^*=\widetilde\varphi u\widetilde\varphi\deleted{\unsure{u^*}}^{\inserted{-1}}$, on a
ici $\Delta(u^*)=\Delta(u)$, donc \deleted{$\unsure{\Delta(u)^2=}$}
\[
\Delta(u)^2=\delta(u)^2
\]
d'où $\Delta(u)=\pm\delta(u)$. Il faut encore examiner
le signe, il faut \deleted{\unsure{pour dire|prendre}} dans un $\operatorname{End}(A)\otimes_{\mathbb Z}\mathbb Q$,
muni de l'involution $u\leadsto u'=\widetilde\varphi u\widetilde\varphi^{-1}$, prouver que
\deleted{$u=u'$}\inserted{$u>0$ (i.e. \deleted{\unsure{vérifiant}} $u=\sum v.v'$)} implique
$\Delta(u)=\uuline{+}\delta(u)$ [N.B. $\delta(u)$ est bien
défini, car $\deg u$ est bien un carré parfait si $u=u'$]

\deleted{Or l'identité $\Delta(u)^2=\delta(u)$ est une identité dans}
% archive sheet number: 24
\scanpage{26}
\deleted{Comme $T_p(u)$ et $T_p(u^*)$ sont}

\deleted{En effet, d'après la théorie p-adique \inserted{de Serre} ou th.
caractéristique de Weil, on aura
\[
(*)\qquad\deg u=\det T_p(u)_{\mathrm{\acute et}}\,\det T_p(u^*)_{\mathrm{\acute et}},
\]
et il reste à prouver que \unsure{$-=$|$=$}
\[
\begin{aligned}
1^\circ)\quad&T_p(u)_{\mathrm{\acute et}}=T_p(u^*)_{\mathrm{\acute et}}\\
2^\circ)\quad&\varepsilon(u)=1
\end{aligned}
\]
\inserted{ce qui implique $\operatorname{d\acute et}(T_p(u))=\varepsilon(u)\sqrt{\deg u}$, $\varepsilon(u)=\pm1$}}

\deleted{D'ailleurs on peut \inserted{se borner à} supposer
$u\in$ $\operatorname{Hom}(A,A)\otimes_{\mathbb Z}\mathbb Q\underset{\varepsilon}{=}\unsure{E|F}$,
et alors $\varepsilon$ est un hom $E^*\to(+1,-1)$, \deleted{\unsure{donc|de plus}}
qui est une fonction \unsure{rationnelle|rationelle} [quotient de
\deleted{\unsure{on voit aussitôt que tout tel hom est|assujettie que tout tel hom est}}
\deleted{identiqumt} deux fonctions}
% archive sheet number: 25
\scanpage{27}
Or pour $u$ variable, $\Delta(u)$ est un polynôme
dans $\unsure{\mathbb Q_p|\mathbb Z_p|\deleted{\mathbb Q}\mathbb Z_p}[T]$ ($T=(T_1\mathbin{-}\mathbin{-}T_{\unsure{N|\alpha}})$, $N=\deleted{\unsure{\dim|\operatorname{rang}}}\operatorname{rg}_{\mathbb Z}\operatorname{End}(A)_{\mathrm{herm}}$)
qui, on vient de le voir, prend des valeurs entières
pour des valeurs entières des variables. Donc
$\Delta(u)$ est un polynôme à coeff. \deleted{\unsure{entiers}} dans $\mathbb Q$
(et non seulemt dans $\mathbb Q_p$), // \deleted{\unsure{donc même à coeff}
dans $\mathbb Z_{(p\mathbb Z)}$ [Car $\mathbb Z_p[T]\cap\mathbb Q[T]=\mathbb Z[T]$,
car $\mathbb Z_p\cap\mathbb Q=\mathbb Z_{(p\mathbb Z)}$]}. Il faut prouver que
ses valeurs sur le cône des $u>0$
\deleted{sont} $\geq0$, et \deleted{pour que} il suffit de le prouver
quand on passe de $\mathbb Q$ à $\mathbb R$. Mais alors
tt \unsure{élmt|élt} $\geq0$ est un carré, d'un élmt
hermitien \inserted{$v$} $>0$, donc $\Delta(u)=\Delta(v^2)=\Delta(v)^2\in\mathbb R^{2}=\mathbb R^+$,
cqfd.

b) $\alpha_p(\varphi)$ induit bien un \deleted{\unsure{hom}}\inserted{isom} de $\operatorname{d\acute et}T_p(A)_{\mathrm{\acute et}}$ avec
$\operatorname{d\acute et}T_p(A^*)_{\mathrm{\acute et}}$. Cela \deleted{\unsure{va déjà pr} \unsure{ces deux|les deux}
(i) $\alpha_p(\varphi)$ applique $\det T_p(A)_{\mathrm{\acute et}}$ sur $\det T_p(A^*)_{\mathrm{\acute et}}$ ;
i.e. $\beta_p(\varphi)$}\inserted{équivaut à la relation}
\[
v_p\bigl(\operatorname{card}(\operatorname{Coker}\deleted{\unsure{\beta_p}}T_p(\widetilde\varphi)_{\mathrm{\acute et}})\bigr)
=v_p\bigl(\sqrt{\deg\widetilde\varphi}\bigr)
\]
Or on a \unsure{évidemment|visiblement}
% archive sheet number: 26
\scanpage{28}
\[
\operatorname{Coker}T_{\deleted{\unsure{p}}p}(\widetilde\varphi)_{\mathrm{\acute et}}\simeq
(\operatorname{Ker}\deleted{T}\widetilde\varphi)(p)_{\mathrm{\acute et}}
\]
Or on sait que $\operatorname{Ker}\widetilde\varphi$ est autodual, donc
$(\operatorname{Ker}\widetilde\varphi)(p)_{\mathrm{\acute et}}$ est dual de $(\operatorname{Ker}\widetilde\varphi(p))_{\mathrm{rad}}$, donc
\unsure{avec|de|comme} $=$ cardinal, d'où la relation
\[
\begin{aligned}
\bigl(\operatorname{card}(\operatorname{Ker}\widetilde\varphi)(p)_{\mathrm{\acute et}}\bigr)^2
&=\operatorname{card}\operatorname{Ker}(\widetilde\varphi(p))\\
&\hspace{2em}\shortparallel\\[-1ex]
&\quad(\deg\widetilde\varphi)(p).
\end{aligned}
\]
d'où la relation envisagée \ldots\ldots\ Cela définit
bien un isom canonique
\[
\boxed{\det T_p(A)_{\mathrm{\acute et}}\xrightarrow[\alpha_A]{\sim}\det T_p(A^*)_{\mathrm{\acute et}}}\tag*{(*)}
\]
compatible avec \unsure{changmt|changt} de base, et isom. en $A$.

[Si $u$ est un autom de $A$, et $u^\vee=u^{*\, -1}$ son
contragrédient, on aura $\alpha\,\det T_p(u)_{\mathrm{\acute et}}=\operatorname{d\acute et}T_p(u^\vee)_{\mathrm{\acute et}}\alpha$,
donc en identifiant \unsure{des|les} dét's \unsure{des|à des} élts de
$\mathbb Z_p$ \unsure{$\times$|ici}, on aura
$\operatorname{d\acute et}T_p(u)_{\mathrm{\acute et}}=\underset{\substack{\shortparallel\\(\operatorname{d\acute et}T_p(u^*))^{-1}}}{\operatorname{det}T_p(u^{*\, -1})_{\mathrm{\acute et}}}$
donc
\[
\boxed{\det T_p(u)_{\mathrm{\acute et}}\inserted{\det}T_p(u^*)_{\mathrm{\acute et}}=1}
\qquad\text{\uline{$u$ un autom de $A$}} \quad]
\]

\margin{\unsure{2|Z}}
Attention, si $u$ est un endom de $A$, on aura
en général $\det T_p(u)\ne\det T_p(u^*)$
% archive sheet number: 27
\scanpage{29}
$=$ si $A$ courbe elliptique sur corps alg. clos ; donc
l'isom $\alpha$ n'est \uline{pas} fonctoriel p.r. à des
hom de schémas abéliens de $=$ dimension
[dans la mesure où ça pourrait avoir un
sens, les deux foncteurs n'étant pas de $=$
variance en $A$ \ldots]. Cependant, \deleted{\unsure{si $u$ est}
\unsure{un hom de $A,B$, on a}}
\[
\boxed{\operatorname{d\acute et} T_p(u)_{\mathrm{\acute et}}\,\operatorname{d\acute et} T_p(u^*)_{\mathrm{\acute et}}=\deg u}
\]
\margin{vaut aussi (via \unsure{identific.|identif.} qu'il permet) pour hom de Sch.A. de $=$ \unsure{dimension|dimensions} rel. $A,B$.}
% slanted

c) Il faudrait vérifier que
\[
\alpha_{A^*}=\alpha_A^{-1}
\]
[sans signe (?)\ldots]. Noter aussi le
comportemt par somme de VA. \ldots

d) L'isom $({}_pA)_{\mathrm{\acute et}}\mathbin{\to}({}_pA^*)_{\mathrm{\acute et}}$ induit par
$\alpha_A$ est bien celui déduit plus haut de
la dualité en cohomologie de DeRham
(à vérifier \ldots).
% archive sheet number: 28
